% Generated draft: requires mathematical and source review.
\documentclass[11pt]{article}
\usepackage[margin=25mm]{geometry}
\usepackage{fontspec}
\setmainfont{Latin Modern Roman}
\usepackage{amsmath,amssymb,mathtools,unicode-math}
\setmathfont{Latin Modern Math}
\usepackage{enumitem,needspace,xurl,hyperref}
\setlength{\parindent}{0pt}
\setlength{\parskip}{5pt}
\setlength{\emergencystretch}{3em}
\newcommand{\N}{\mathbb{N}}
\newcommand{\Z}{\mathbb{Z}}
\newcommand{\Q}{\mathbb{Q}}
\newcommand{\R}{\mathbb{R}}
\newcommand{\C}{\mathbb{C}}
\newcommand{\F}{\mathbb{F}}
\newcommand{\A}{\mathbb{A}}
\newcommand{\PP}{\mathbb{P}}
\newcommand{\E}{\mathbb{E}}
\newcommand{\eps}{\varepsilon}
\newcommand{\dd}{\,\mathrm{d}}
\begin{document}
\section*{UnsolvedMath ID: 30003356}
\textbf{Source problem code:} catalogue entry
\subsection{Definitions and mathematical statement}
Let \(\mathbb{CP}^{3}\) denote the space of complex lines in \(\C^{4}\), and let \(\pi\colon M\to S^{1}\) be a smooth fiber bundle with fiber \(\mathbb{CP}^{3}\). For a diffeomorphism \(\tau\colon\mathbb{CP}^{3}\to\mathbb{CP}^{3}\), define its mapping torus by
\[
T_{\tau}=\bigl(\mathbb{CP}^{3}\times[0,1]\bigr)\big/\bigl((x,1)\sim(\tau(x),0)\bigr).
\]
This explains the notation \(S^{1}\times_{\tau}\mathbb{CP}^{3}\) used in the supplied statement. However, the supplied record does not define its intended \(\tau\). The questions below therefore retain an unspecified intended diffeomorphism \(\tau\); they do not quantify universally or existentially over \(\tau\).

A smoothly varying family of metrics on the fibers means a smooth positive definite symmetric bilinear form \(h\) on the vertical tangent bundle \(\ker(\dd\pi)\). It induces a Riemannian metric \(h_b\) on each fiber \(M_b=\pi^{-1}(b)\). Here the phrase ``positively curved'' is interpreted as positive sectional curvature on each fiber; this interpretation requires confirmation against the source.

For every such smooth bundle, the supplied record asks the following two separate questions:
\begin{enumerate}
\item If there exists a Riemannian metric \(g\) on \(M\) such that \(\sec_g(\sigma)\geq 0\) for every point \(p\in M\) and every two-dimensional subspace \(\sigma\subset T_pM\), must
\[
M\cong S^{1}\times\mathbb{CP}^{3}\quad\text{or}\quad M\cong T_{\tau}
\]
hold, where \(\cong\) denotes diffeomorphism of total spaces?
\item If there exists a smooth vertical metric \(h\) such that \(\sec_{h_b}(\sigma)>0\) for every \(b\in S^{1}\), every \(p\in M_b\), and every two-dimensional subspace \(\sigma\subset T_pM_b\), must the same diffeomorphism conclusion hold?
\end{enumerate}
The second question is stated independently of the first curvature hypothesis. The conclusion as supplied concerns diffeomorphism, without requiring preservation of the bundle projection. Identification of the intended \(\tau\) is necessary to make the proposed classification fully specified.
\subsection{Short English statement}
For a smooth \(\mathbb{CP}^{3}\)-bundle over the circle, ask separately whether nonnegative sectional curvature on its total space, or a smooth family of positively sectionally curved fiber metrics, forces its total space to be diffeomorphic to the product or the specified twisted mapping torus. The record leaves the twisting diffeomorphism \(\tau\) undefined.
\subsection{Sources}
\begin{itemize}
\item Supplied local UnsolvedMath registry record, permanent ID \(30003356\), source problem code \texttt{catalogue entry}, title ``Curvature Rigidity of Complex Projective 3-Space Bundles,'' in the Oberwolfach Reports Open Problems collection. No individual proposer is supplied; the proposed year is recorded as 2017.
\item Supplied source citation: \emph{Mini-Workshop: Spaces and Moduli Spaces of Riemannian Metrics} (2018), \emph{Oberwolfach Reports}, pp.\ 133--166, DOI \texttt{10.4171/owr/2017/3}. Also listed as ``Oberwolfach report on CP3 bundles'': \url{https://doi.org/10.4171/owr/2017/3}.
\item Supplied literature reference: ``Moduli spaces of positive curvature metrics in dimension four and beyond,'' \url{https://doi.org/10.1007/s00209-026-03947-3}.
\item Literature assessment dated exactly \texttt{2026-08-21}: the supplied curation labels the problem open, describes both classification questions as remaining broad diffeomorphism and curvature-rigidity problems, and states that the cited report gives no theorem resolving either alternative. This is supplied dated curation, not an independently verified literature assessment.
\item The registry labels the statement \texttt{corrected\_verified} and reports that the official report passage was checked and rewritten while preserving its mathematical content. These are supplied verification claims; no independent source inspection is asserted here.
\end{itemize}
\end{document}
